\itemhead{1.2 例题 · 例题5（教材p.18）}{XJ-B1-C01-EX-P018-Y0491}
\label{XJ-B1-C01-EX-P018-Y0491}
{\footnotesize 湘教版教材p.18 / PDF第25页 / 1.2 例题 / 例题5}
\field{原题}{试证：
\subq{1}在实数范围内，$x=1$ 是 $x^2=1$ 的充分而不必要条件；
\subq{2}四边形的两组对边分别相等是四边形为矩形的必要而不充分条件；
\subq{3}在 $\triangle ABC$ 中，$\angle B=\angle C$ 是 $AC=AB$ 的充要条件。}
\field{参考答案}{（1）$x=1\Rightarrow x^2=1$，反向有反例 $x=-1$；（2）矩形对边分别相等，斜的平行四边形说明反向不成立；（3）利用等边对等角和等角对等边双向推出。}
{\small \field{关联知识}{L02 充分、必要、充要条件（教材帮p.30、31 / PDF39、40）}
\field{方法与易错点}{分别检查两个方向；用解集包含或具体反例判定，不凭句子位置猜方向。 易错：把“必要”箭头反写；只有一个方向证据便断言“不必要”；忽略给定论域。}
\field{答案依据}{例题原解位置：湘教版教材p.18及题库原图；独立编写可编辑参考解答。答案册此部分主要收练习、习题，不伪造例题答案册位置。}}
